% ABSTRACT
% Replace this with your actual abstract content

This thesis examines the predictive power of convolutional neural networks (CNNs) for stock returns, with a particular focus on Federal Open Market Committee (FOMC) announcement days. I replicate the CNN-based prediction methodology of Jiang et al. (2023) and extend it to analyze cross-sectional return predictability during high-information macro events. Using a sample of 22,480 U.S. stocks from 2001-2024, I find that CNN predictions generate economically significant returns: equal-weighted high-minus-low portfolios earn 70.74\% annually with a Sharpe ratio of 5.60, while value-weighted portfolios earn 22.69\% annually. 

The main contribution is an event study of 217 FOMC meetings, which reveals significant cross-sectional predictability on announcement days (0.21\%, t=2.95, p$<$0.01) and in intermediate windows (0.35\%, t=2.24, p$<$0.05). However, when compared to matched non-FOMC days, FOMC periods show uniformly lower predictability across all horizons (differences of -0.49\% to -0.87\%, all t$>$4, p$<$0.001). This pattern supports the hypothesis that heightened market attention during macro events increases efficiency and compresses exploitable cross-sectional patterns, consistent with Hirshleifer and Sheng (2021). 

The predictability is concentrated in small-cap stocks (equal-weighted returns are 3-10× larger than value-weighted), suggesting that limited attention and limited arbitrage are key mechanisms. These findings demonstrate that machine learning predictability is state-dependent and varies with the information environment, providing new evidence on the interaction between attention, efficiency, and cross-sectional return predictability.

